\section{Reproducible computation and the meaning of a certificate}
\label{sec:computation}

The archive is organized so that the analytic argument can be read
without executing code, while each finite algebraic step can be replayed
without trusting a decimal fit. There are three distinct kinds of
computation. Exact rational word algebra verifies a shuffle certificate.
Outward interval arithmetic, together with a proved analytic tail bound,
certifies a numerical enclosure. Ordinary high-precision quadrature checks
an implementation and its conventions. These roles should remain
distinguishable when the files are integrated into ProveIt.

\subsection{Exact word algebra and algebraic substitutions}

The program \code{code/shuffle_certificates.py} represents a binary word
as a tuple and a shuffle as a dictionary of integer multiplicities.
It recursively shuffles two words, computes the nonincreasing Lyndon
factorization, and solves the resulting triangular equation using
rational numbers. The decisive replay is to expand the resulting
polynomial back into words. A successful replay returns exactly the
original word with coefficient one and no other terms.

At weights one through nine, this procedure re-expands all \(511\)
nonempty words ending in one. It also enumerates all \(45\) nonempty
bigraded sectors and checks their Lyndon counts against
Theorem~\ref{thm:witt}. The recorded data contain every polynomial normal
form of the ten weight-six triples, all four weight-five doubles, and
the two-product certificate of Corollary~\ref{cor:wt5}.
This exhaustive finite check supports the implementation; the
all-weight statement rests on the proof of the theorem.

The depth verifier implements the binomial formula independently of the
target expressions. Exact symbolic subtraction confirms every coefficient
in the five weight-six and three weight-four source formulas.
It also emits the parity-reducible Gaussian component of every double
through weight twelve: imaginary parts at even weight and real parts at
odd weight. Thus the archive includes further explicit identities as
reusable data, with no integer-relation recognition step.

The ladder verifier works in
\(\Q[r]/(r^3+r-1)\). Rational functions are reduced using the inverse
of their denominator modulo the defining polynomial. Its \(35\) exact
checks cover the six factorizations, the eighteen Kummer arguments,
the eight Rogers arguments, the rational row operation, and both
logarithmic-tail cancellations. In particular, the finite certificate
does not attempt to evaluate a transcendental function in the number
field: it verifies the arguments and coefficients of functional equations
whose analytic validity is stated in the article.

\subsection{Independent numerical diagnostics}

For the same-argument double family, a useful independent evaluator is
\begin{equation}
F_{a,b}(z)=\frac{z}{(a-1)!}
\int_0^1\frac{(-\log t)^{a-1}\Li_b(zt)}{1-zt}\,dt,
\qquad |z|=1,\ z\ne1.
\label{eq:double-quadrature}
\end{equation}
It follows first by absolutely convergent series expansion for \(|z|<1\),
then by the boundary limit. This evaluator uses neither the binomial
reduction nor a fitted constant basket.
For the inverse-argument mixed family the corresponding integral replaces
\(\Li_b(zt)\) by \(\Li_b(t)\).
For the arbitrary-depth one-two family the verifier uses the integral
with the zero letter held fixed, before expanding its outer factor.

The resulting diagnostics comprise \(64\) checks at \(65\) decimal
working digits for the same-argument double and one-two formulas,
including both Gaussian and Eisenstein points and the non-cyclotomic
point \(e^{i\sqrt2}\). Their largest absolute residual was
\(1.081\cdot10^{-64}\) when rounded upward.
A separate inherited-mixed audit contains eighteen general-formula
quadratures and four explicit specialization checks, with maximum
residual less than \(1.365\cdot10^{-65}\).
The cubic and complementary-power ladders were evaluated at \(100\)
and \(200\) working digits, including comparison with directly summed
convergent classical-polylogarithm series. The cubic residual at
\(200\) digits was below \(2.62\cdot10^{-200}\); the independent
series comparison differed by less than \(9.91\cdot10^{-191}\).

For the remaining \(S_4\) conjecture, an independent real integral is
\[
S_4=-\frac16\int_0^1
\frac{(-\log t)^3\log(1+t^2)}{1+t^2}\,dt.
\]
It gives
\[
S_4=-0.0104934562973350434533567510062017519635106188773\ldots.
\]
At \(50\) working digits the discrepancy from
\eqref{eq:S4-short} was below \(5.0\cdot10^{-51}\).
This is evidence for the conjecture, and is deliberately excluded from
the proved-identity count. The archive includes a small independent
script for this comparison.

None of these residuals is an interval-certified error bound. In
particular, a residual at approximately the working precision is
consistent with cancellation and finite-precision rounding.
The independent formulas and the analytic proofs, rather than the
number of matching digits, determine the mathematical status.

\subsection{The interval certificate and the asymptotic diagnostics}

The cubic moment calculation has a stronger status. All endpoints in
\code{code/moments/cubic_certificate.json} are exact rational numbers
with denominator \(10^{130}\), displayed in decimal form. The
certificate uses no floating-point special-function values. The interval
width is less than \(2.387\cdot10^{-113}\), and its endpoints share
\(112\) digits after the decimal point. This certifies the common
truncated decimal prefix; a claim about rounded digits would instead
require checking the relevant rounding boundary.

The certificate was replayed independently and reproduced every recorded
field exactly. Separate quadratures at \(150\) and \(180\) working
digits lie inside the interval. The quadratures provide a useful
cross-check but are not used to construct the rational enclosure.

The coefficient asymptotic and moment bounds were also compared with
direct evaluations. Figure~\ref{fig:moment-asymptotics} shows two
features that matter for using the expansion. The coefficient
approximation improves in the predicted inverse powers of \(k\).
Nevertheless, at each fixed \(n\), the terms \(c_k/k^n\) eventually
increase. It is therefore essential to specify a truncation rule,
such as the one in Theorem~\ref{thm:effective-moment}.

\begin{figure}[htbp]
\centering
\includegraphics[width=\linewidth]{figures/moment_asymptotics.pdf}
\caption{Numerical diagnostics for the proved coefficient and moment
asymptotics. Left: absolute relative errors of
\(C\rho^{-k}k^{-3/2}\) and
\(C\rho^{-k}k^{-3/2}(1+\beta_1/k)\), using coefficients computed from
their positive defining recurrence. Right:
\(\log_{10}(c_k/k^n)\) for \(n=20,40,80\). Dotted vertical lines
mark \(k=n/\alpha\); circles mark the sampled minima.
All plotted quantities are dimensionless. The figure illustrates
the theorems and supplies no independent error certificate.}
\label{fig:moment-asymptotics}
\end{figure}

\subsection{Archive structure and reproduction}

The top-level \code{README.md} gives installation, build, verification,
and integration instructions. The article's master source is
\code{article/polylogarithms_exact_reductions.tex}; all included
sections and the vector figure are supplied. The verification programs
are grouped under \code{code/}, with their recorded outputs alongside
the relevant program or under \code{data/}. Source provenance includes
the pinned commit, content hashes, and a status register mapping each
proposed manuscript update to its source equation or section.

From the archive root, the essential commands are
\begin{quote}\small\ttfamily
python -m pip install -r requirements.txt\\
python code/run\_checks.py --fast\\
python code/run\_checks.py --full\\
make pdf
\end{quote}
The fast command replays exact algebra and the rational interval
certificate. The full command additionally runs the independent
quadratures and regenerates the diagnostic figure. It may take several
minutes; no network access is required after installing the dependencies.
The complete recorded diagnostics remain available without rerunning
them. The PDF is built with ordinary \LaTeX{} packages and contains
no external data dependency.

The source manifest is a checkpoint, not an assertion that later
repository revisions have the same status. Integration should preserve
the proved/conjectured distinction, the descending-index convention,
the branch choices, and the qualification that shuffle quotient
dimensions are formal. The supplied register is intended to make
that editorial task precise.
